%==========================================================================
%  TRES0312 Fysikk — Foiler-mal
%
%  MØrk/lys modus: kommenter ut én av de to linjene nedenfor.
%==========================================================================
\newif\ifdarkmode
\darkmodetrue      % ← mørk modus
%\darkmodefalse   % ← lys modus  (fjern % foran og legg % foran linjen over)

\documentclass[aspectratio=169, 12pt]{beamer}

\usepackage{amd-beamer}
\usetikzlibrary{overlay-beamer-styles}   % for "visible on=<->" i enkelte tikz-figurer

%--------------------------------------------------------------------------
%  DOKUMENTINFO
%--------------------------------------------------------------------------
\title{Kinematikk}
\subtitle{TRES0312 --- Sirkelbevegelse}
\author{Ben David Normann}
\date{\today}

%--------------------------------------------------------------------------
%  SEKSJONSKONTEKST
%  \setcounter{section}{N} setter gjeldende seksjonsnummer til N, slik at
%  \defn, \exm og \oppgave får samme nummerering som i kompendiet.
%  Tellerne nullstilles automatisk ved hvert \setcounter{section}{...}.
%--------------------------------------------------------------------------
\setcounter{section}{3}   % Kapittel 3: "Kinematikk"

\begin{document}

%--------------------------------------------------------------------------
%  SLIDE 1 — TITTELSIDE
%--------------------------------------------------------------------------
\begin{frame}[plain]
  \titlepage
  \dekkerdelkap{kapittel 3.5}
\end{frame}

%--------------------------------------------------------------------------
%  SLIDE 2 — Til ettertanke: retningen til akselerasjonen i sirkelbevegelse
%--------------------------------------------------------------------------
\begin{frame}{}

  \tanke{}{%
  Hvilken retning peker akselerasjonen i sirkelbevegelse?
  }

  \pause
  \begin{figure}[h]
  \centering
  % LEFT: position diagram
  \begin{minipage}[t]{0.50\textwidth}
  \centering
  \begin{tikzpicture}[font=\small, scale=0.95, >=stealth]
    % Centre O
    \fill (0,0) circle (2pt);
    \node[below left=2pt] at (0,0) {$O$};

    % Partial circle arc (40° – 140°)
    \draw[headCol, thick]
      ({2*cos(40)},{2*sin(40)}) arc[start angle=40, end angle=140, radius=2];

    % Points P1 (110°) and P2 (70°)
    \coordinate (P1) at ({2*cos(110)},{2*sin(110)});
    \coordinate (P2) at ({2*cos(70)},{2*sin(70)});

    % Radii
    \draw[black!40, semithick] (0,0) -- (P1);
    \draw[black!40, semithick] (0,0) -- (P2);

    % r labels at midpoints of radii
    \node[left=2pt,  black!55, font=\footnotesize] at ({cos(110)},{sin(110)}) {$r$};
    \node[right=2pt, black!55, font=\footnotesize] at ({cos(70)}, {sin(70)})  {$r$};

    % Angle Δθ arc at O
    \draw ({0.52*cos(110)},{0.52*sin(110)})
          arc[start angle=110, end angle=70, radius=0.52];
    \node[above=1pt, font=\footnotesize] at (0, 0.62) {$\Delta\theta$};

    % Points
    \fill[defscol] (P1) circle (2.5pt);
    \node[defscol, above left=2pt] at (P1) {$P_1$};
    \fill[defscol] (P2) circle (2.5pt);
    \node[defscol, above right=2pt] at (P2) {$P_2$};

    % Chord Δs (P1 → P2, horizontal since symmetric about 90°)
    \draw[->, oppsumcol, thick, dashed] (P1) -- (P2)
          node[midway, below=5pt] {$\Delta s$};

    % Velocity v1 at P1 (tangent = 110° − 90° = 20°)
    \draw[->, very thick, defscol]
          (P1) -- ++({cos(20)},{sin(20)})
          node[above right=-1pt] {$\vec{v}_1$};

    % Velocity v2 at P2 (tangent = 70° − 90° = −20°)
    \draw[->, very thick, defscol]
          (P2) -- ++({cos(-20)},{sin(-20)})
          node[right=2pt] {$\vec{v}_2$};
  \end{tikzpicture}
  \end{minipage}%
  \hfill
  % RIGHT: velocity triangle
  \begin{minipage}[t]{0.46\textwidth}
  \centering
  \begin{tikzpicture}[font=\small, scale=0.95, >=stealth]
    % Tails at origin, v_len = 1.5
    \coordinate (tail) at (0,0);
    \coordinate (tip1) at ({1.5*cos(20)}, { 1.5*sin(20)});
    \coordinate (tip2) at ({1.5*cos(20)}, {-1.5*sin(20)});

    % v1 and v2
    \draw[->, very thick, defscol] (tail) -- (tip1);
    \draw[->, very thick, defscol] (tail) -- (tip2);

    % v labels along the arrows
    \node[defscol, above left=1pt, font=\footnotesize]
          at ({0.75*cos(20)},{ 0.75*sin(20)}) {$v$};
    \node[defscol, below left=1pt, font=\footnotesize]
          at ({0.75*cos(20)},{-0.75*sin(20)}) {$v$};

    % Arrow-head labels (v1, v2)
    \node[defscol, above=2pt]  at (tip1) {$\vec{v}_1$};
    \node[defscol, below=2pt]  at (tip2) {$\vec{v}_2$};

    % Δv from tip1 to tip2 (straight down)
    \draw[->, very thick, oppsumcol] (tip1) -- (tip2)
          node[right=3pt, midway] {$\Delta\vec{v}$};

    % Tail point
    \fill (tail) circle (1.5pt);

    % Angle Δθ at tail
    \draw ({0.52*cos(20)},{0.52*sin(20)})
          arc[start angle=20, end angle=-20, radius=0.52];
    \node[right=3pt, font=\footnotesize] at (0.60, 0) {$\Delta\theta$};

    % Bounding-box adjustment (match height of left panel)
    \path (0,0); \node[opacity=0] at (0,2.5) {$x$};
    \path (0,-1.7);
  \end{tikzpicture}
  \end{minipage}
  \end{figure}

\end{frame}

%--------------------------------------------------------------------------
%  SLIDE 3 — Definisjon: Sentripetalakselerasjon
%--------------------------------------------------------------------------
\begin{frame}{Sentripetalakselerasjon}

  \defnat{3.6}{Sentripetalakselerasjon}{
    Sentripetalakselerasjonen er akselerasjonen inn mot sentrum i ren sirkelbevegelse med konstant fart. Uttrykket for størrelsen på denne akselerasjonen kan vises å være
    \[
      a_\perp = \frac{v^2}{r}
    \]
    hvor $v$ er banehastigheten til objektet og $r$ er radiusen til sirkelen.
  }

  \pause
  \merk{Retning}{
  Sentripetalakselerasjonen peker \emph{alltid} mot sentrum av sirkelen --- selv om objektet hele tiden beveger seg langs kanten av sirkelen.
  }

\end{frame}

%--------------------------------------------------------------------------
%  SLIDE 4 — Aktivitet: Bil i rundkjøring
%--------------------------------------------------------------------------
\begin{frame}{Aktivitet}

  \aktivitet{Bil i rundkjøring}{%
  En bil kjører gjennom en rundkjøring med konstant fart $v = 36\,\mathrm{km/h}$ og radius $r = 20\,\mathrm{m}$.
  Hva er sentripetalakselerasjonen til bilen?
  }

\end{frame}

%--------------------------------------------------------------------------
%  SLIDE 5 — Aktivitet: Månen i bane rundt jorda
%--------------------------------------------------------------------------
\begin{frame}{Aktivitet}

  \aktivitet{Sirkelbevegelse og omløpstid}{%
  Månens gjennomsnittlige avstand fra jorda er $r = 3.84\cdot10^8\,\mathrm{m}$, og omløpstiden er $T = 27.3\,\text{dager}$.
  \begin{enumerate}[a)]
    \item Gjør om omløpstiden til sekunder.
    \item Finn sentripetalakselerasjonen månen opplever ved hjelp av formelen $a_\perp = 4\pi^2 r / T^2$.
    \item Beregn månens banehastighet $v = 2\pi r / T$ og verifiser svaret ved å bruke $a_\perp = v^2/r$.
  \end{enumerate}
  }

\end{frame}

%--------------------------------------------------------------------------
%  SLIDE 6 — Oppsummering: Omløpstid
%--------------------------------------------------------------------------
\begin{frame}{}

  \oppsum{Omløpstid}{%
  Farten i sirkelbevegelse kan skrives ved hjelp av omløpstiden $T$ (tiden det tar å fullføre ett omløp):
  \[
    v = \frac{2\pi r}{T}.
  \]
  Setter vi dette inn i uttrykket for sentripetalakselerasjon, $a_\perp = v^2/r$, får vi et alternativt uttrykk for akselerasjonen (likning 3.16 i kompendiet):
  }

  \graabox{$\displaystyle a_\perp = \frac{4\pi^2 r}{T^2}$}

\end{frame}

%--------------------------------------------------------------------------
%  SLIDE 7 — Neste gang
%--------------------------------------------------------------------------
\begin{frame}
\begin{center}
  \Huge{Neste gang: Dynamikk --- \textit{``Hvorfor bevegelse!''}}
  \begin{itemize}
    \item Newtons lover
    \item Kraft-begrepet
    \item Krefter i dagliglivet (tyngdekraft, normalkraft, friksjon)
  \end{itemize}
  \vspace{6mm}
  {\normalsize\itshape May the force be with you\ldots}
\end{center}
\end{frame}

\end{document}
